\documentclass[]{article}
\usepackage{geometry,tikz,lipsum,amsmath,amssymb,soul,xcolor,cancel}
\tikzset{every picture/.style={line width=0.75pt}} %set default line width to 0.75pt        
\numberwithin{equation}{section}

\NewDocumentCommand{\caja}{ m O{gray!40} O{gray!10} m }{
	\begin{figure}[h!]
		\centering
		\begin{tikzpicture}
			\node[anchor=text, text width=\columnwidth-1.2cm, draw, rounded corners, line width=1pt, fill=#3, inner sep=5mm] (big) {\\#4};
			\node[draw, rounded corners, line width=.5pt, fill=#2, anchor=west, xshift=5mm] (small) at (big.north west) {#1};
		\end{tikzpicture}
	\end{figure}
}

\newcommand{\deriv}[3]{\left(\frac{\partial #1}{\partial #2} \right)_{#3}}

\newcommand{\dderiv}[3]{\left(\dfrac{\partial #1}{\partial #2} \right)_{#3}}

\NewDocumentCommand{\indUpDown}{ O{\Lambda} O{\mu} O{\nu} }{
	{#1}^{#2}_{\ #3}
}
\NewDocumentCommand{\indDownUp}{ O{\Lambda} O{\mu} O{\nu} }{
	{#1}_{#2}^{\ #3}
}

\newcommand{\LambdaMuNu}{\indUpDown}
\newcommand{\0}{\textcolor{gray}{0}}
\newcommand{\ei}{\mathbf{e}_i}
\newcommand{\ej}{\mathbf{e}_j}
\newcommand{\lc}{\epsilon_{ijk}}
\newcommand{\gr}{\nabla}
\newcommand{\di}{\nabla\cdot}
\newcommand{\ro}{\nabla\times}
\newcommand{\la}{\nabla^2}
\newcommand{\rv}{\mathbf{r}}
\newcommand{\pv}{\mathbf{p}}
\newcommand{\xv}{\mathbf{x}}
\newcommand{\yv}{\mathbf{y}}
\newcommand{\zv}{\mathbf{z}}
\newcommand{\xhat}{\hat{\xv}}
\newcommand{\yhat}{\hat{\yv}}
\newcommand{\zhat}{\hat{\zv}}
\newcommand{\ihat}{\hat{\mathbf{i}}}
\newcommand{\jhat}{\hat{\mathbf{j}}}
\newcommand{\khat}{\hat{\mathbf{k}}}
\newcommand{\rhat}{\hat{\rv}}
\newcommand{\that}{\hat{\mathbf{\theta}}}
\newcommand{\phat}{\hat{\mathbf{\varphi}}}
\newcommand{\rhhat}{\hat{\mathbf{\rho}}}
\newcommand{\nhat}{\hat{\mathbf{n}}}
\newcommand{\Rhat}{\hat{\mathbf{R}}}
\newcommand{\eo}{\epsilon_{0}}
\newcommand{\uo}{\mu_{0}}
\NewDocumentCommand{\ke}{ O{1}}{
	\frac{#1}{4\pi\eo}
}
\NewDocumentCommand{\ku}{ O{}}{
	\frac{\uo #1}{4\pi}
}
%\newcommand{\ke}{\frac{1}{4\pi\eo}}
%\newcommand{\ku}{\frac{\uo}{4\pi}}
\newcommand{\rdif}{\rv - \rv'}
\newcommand{\pd}[2]{\dfrac{\partial #1}{\partial #2}}
\newcommand{\lpd}[2]{\partial #1/\partial #2}
\newcommand{\td}[2]{\dfrac{d #1}{d #2}}
\newcommand{\ltd}[2]{d #1/d #2}
\newcommand{\n}[1]{\mathbf{#1}}
\newcommand{\eff}[1]{#1_\mathrm{eff}}
\newcommand{\Ueff}{\eff{U}}

\newcommand{\wfrac}[3][3pt]{%
	\frac{\wfracterm{depth}{\dp}{#1}{#2}}{\wfracterm{height}{\ht}{#1}{#3}}%
}
\newcommand{\wfracterm}[4]{%
	\sbox0{$\displaystyle#4$}%
	\vrule width 0pt #1 \dimexpr #20+#3\relax
	\usebox{0}%
}

\newcommand{\pr}[1]{\left\langle #1 \right\rangle }

\newcommand{\qed}{\hfill\(\blacksquare\)}

%opening
\title{Certamen 1 - Física Estadística Avanzada}
\author{Simon Fonseca}

\begin{document}
\maketitle

\section{Problema 1}
Considere un sistema clásico de moléculas diatómicas no interactuantes encerradas en un a caja de volumen $V$ a temperatura $T $. El Hamiltoniano de una única molécula está dado por:
\begin{equation}
	H(\rv_1,\rv_2,\pv_1,\pv_2) = \frac{1}{2m}\left(p_1^2 + p_2^2\right) + \frac{1}{2} K \left(\rv_1 - \rv_2\right)^2
\end{equation}
Estudie la termodinámica de este sistema, incluyendo la dependencia de la cantidad $\pr{r_{12}^2}$ a temperatura $T $.
\\

\textbf{Respuesta:}
Tenemos que los valores fijos del sistema son el volumen $V$, la temperatura $T$ y el número de partículas $N$, por lo que usaremos el ensamble canónico para describir al sistema:
\begin{align*}
	Z(T,V,N) &= \frac{1}{N! h^{3N}}\int e^{-\beta H(\Gamma)} \, d\Gamma \nonumber\\
	&= \frac{1}{N! h^{6N}} \int e^{-\beta H(\rv_1,\rv_2,\pv_1,\pv_2)} \, d^3\rv_{1,1} d^3\rv_{2,1} \ldots d^3\rv_{1,N} d^3\rv_{2,N} \, d^3\pv_{1,1} d^3\pv_{2,1} \ldots d^3\pv_{1,N} d^3\pv_{2,N} \\
	&= \frac{1}{N! h^{6N}} \left[\int e^{-\beta H(\rv_1,\rv_2,\pv_1,\pv_2)} \, d^3\rv_1 d^3\rv_2 \, d^3\pv_1 d^3\pv_2\right]^N \\
	&= \frac{1}{N! h^{6N}} \left[\int e^{-\beta \frac{1}{2m}\left(p_1^2 + p_2^2\right)} \, e^{-\beta \frac{1}{2} K \left(\rv_1 - \rv_2\right)^2} \, d^3\rv_1 d^3\rv_2 \, d^3\pv_1 d^3\pv_2\right]^N \\
	&= \frac{1}{N! h^{6N}} \left[\int e^{-\beta \frac{1}{2m}\left(p_1^2 + p_2^2\right)} \, d^3\pv_1 d^3\pv_2 \, \int e^{-\beta \frac{1}{2} K \left(\rv_1 - \rv_2\right)^2} \, d^3\rv_1 d^3\rv_2 \right]^N 
	\intertext{Pero $e^{p_1^2} = e^{p_{1x}^2 + p_{1y}^2 + p_{1z}^2} = e^{p_{1x}^2} + e^{p_{1y}^2} + e^{p_{1z}^2}$, por lo que}
	&= \frac{1}{N! h^{6N}} \left[\left(\int e^{-\beta \frac{1}{2m} p^2} \, dp \right)^6 \, \int e^{-\beta \frac{1}{2} K \left(\rv_1 - \rv_2\right)^2} \, d^3\rv_1 d^3\rv_2 \right]^N 
\end{align*}
Pero, considerando el cambio de variables \(\rv_1 = \n{R} + \rv/2, ~~ \rv_2 = \n{R} - \rv/2\)
\begin{align*}
	e^{\left(\rv_1 - \rv_2\right)^2} d^3\rv_1 d^3\rv_2 &= e^{\rv^2} \left|\pd{(\rv_1,\rv_2)}{(\n{R},\rv)}\right| \, d^3\n{R} d^3\rv\\
	&= e^{\rv^2} |-1| \, d^3\n{R} d^3\rv
%	e^{\left(r_{1x} - r_{2x}\right)^2 + \left(r_{1y} - r_{2y}\right)^2 + \left(r_{1z} - r_{2z}\right)^2}  \, dr_{1x} dr_{1y} dr_{1z} dr_{2x} dr_{2y} dr_{2z}\\
%	&= \left(e^{\left(r_{1} - r_{2}\right)^2} \, dr_1 dr_2\right)^3
%	&= \left(e^{r_{1}^2 - 2 r_{1}r_{2} + r_{2}^2} \, dr_1 dr_2\right)^3\\
%	&= \left(e^{r_{1}^2} e^{- 2 r_{1}r_{2}} e^{r_{2}^2} \, dr_1 dr_2\right)^3\\
%	&= \left(e^{r^2} \, dr \right)^6 \left( e^{- r_{1}r_{2}} \, dr_1 dr_2\right)^6
\end{align*}
Integrando con \texttt{Mathematica} en todo el espacio,
\begin{align}
	Z(T,V,N) &= \frac{1}{N! h^{6N}} \left[\left(\int e^{-\beta \frac{1}{2m} p^2} \, dp \right)^6 \, \int \, d^3\n{R} \, \int e^{-\beta \frac{1}{2} K \rv^2} \, d^3\rv \right]^N \nonumber\\
	&= \frac{1}{N! h^{6N}} \left[\left( \sqrt{\frac{2\pi m}{\beta}} \right)^6 \, V \, \left(\sqrt{\frac{2\pi}{\beta K}}\right)^3 \right]^N \nonumber\\
	&= \frac{1}{N! h^{6N}} \left[V \left( \frac{8\pi^3 m^2}{\beta^3 K} \right)^{3/2} \right]^N 
\end{align}

Ahora que tenemos una expresión para la función de partición, podemos utilizarla para calcular la energía libre de Helmholtz,
\begin{align*}
	F &= -\frac{1}{\beta} \ln Z \\
	&= -\frac{1}{\beta}\left[-\ln\left(N! h^{6N}\right) + N \ln\left(V\right) + \frac{3N}{2}\ln\left(\frac{8\pi^3 m^2}{\beta^3 K}\right) \right]
\end{align*}
Con la aproximación de Stirling,
\begin{align}
	\ln Z &= -\ln\left(N!\right) - \ln\left(h^{6N}\right) + N \ln\left(V\right) + \frac{3N}{2}\ln\left(\frac{8\pi^3 m^2}{\beta^3 K}\right) \nonumber\\
	\ln Z &\approx N - N\ln N - 6N\ln h + N \ln V + \frac{3N}{2} \ln\left(\frac{8\pi^3 m^2}{K}\right) - \frac{9N}{2} \ln \beta \\
	\Rightarrow F &\approx -\frac{1}{\beta}\left[N - N\ln N - 6N\ln h + N \ln V + \frac{3N}{2}\ln\left(\frac{8\pi^3 m^2}{K}\right) - \frac{9N}{2}\ln \beta \right]
\end{align}

Y ahora podemos obtener varias cantidades termodinámicas derivadas de la expresión de Helmholtz,
\begin{align}
	U &= - \pd{\ln Z}{\beta} \approx \frac{9N}{2\beta}\\
	p &= \frac{1}{\beta} \pd{\ln Z}{V} \approx \frac{N}{V \beta}
\end{align}
Y la entropía,
\begin{equation}
	S = k_B\beta(U - F) \approx N k_B\left[\frac{11}{2} - 6\ln h + \ln N - \ln V - \frac{3}{2}\ln\left(\frac{8\pi^3 m^2}{K}\right) + \frac{9}{2} \ln \beta \right]
\end{equation}

Finalmente, la cantidad $\pr{r_{12}^2}$, es el valor promedio de la distancia entre los dos átomos de una molécula diatómica. Para calcularlo, consideramos que \(r_{12}^2 = |\rv_1 - \rv_2|^2 = \left(\rv_1 - \rv_2\right)^2 \), tendremos, para una sola molécula:
\begin{align*}
	\pr{r_{12}^2} &= \frac{\int \left(\rv_1 - \rv_2\right)^2 e^{-\beta \Gamma} \, d\Gamma}{\int e^{-\beta \Gamma} \, d\Gamma} \\
	&= \left(\int e^{-\beta \frac{1}{2m} p^2} \, dp \right)^6 \, \int \left(\rv_1 - \rv_2\right)^2 e^{-\beta \frac{1}{2} K \left(\rv_1 - \rv_2\right)^2} \, d^3\rv_1 d^3\rv_2  \left[V \left( \frac{8\pi^3 m^2}{\beta^3 K} \right)^{3/2} \right]^{-1}\\
	&= \left( \frac{2\pi m}{\beta} \right)^3 \, \int \left(\rv_1 - \rv_2\right)^2 e^{-\beta \frac{1}{2} K \left(\rv_1 - \rv_2\right)^2} \, d^3\rv_1 d^3\rv_2 ~ \frac{1}{V} \left( \frac{\beta^3 K}{8\pi^3 m^2} \right)^{3/2}
	\intertext{Realizando el mismo cambio de variables, \(\rv_1 = \n{R} + \rv/2, ~~ \rv_2 = \n{R} - \rv/2\),}
	&= \left( \frac{2\pi m}{\beta} \right)^3 \, \int \, d^3\n{R} \, \int \rv^2 e^{-\beta \frac{1}{2} K \rv^2} \, d^3\rv ~ \frac{1}{V} \left( \frac{\beta^3 K}{8\pi^3 m^2} \right)^{3/2}\\
	&= \int \rv^2 e^{-\beta \frac{1}{2} K \rv^2} \, d^3\rv ~ \left( \frac{\beta^3 K}{8\pi^3 m^2} \right)^{3/2}\\
	&= \frac{6\sqrt{2} \pi^{3/2}}{(\beta K)^{5/2}} ~ \left( \frac{\beta^3 K}{8\pi^3 m^2} \right)^{3/2}\\
	&= \frac{3}{\beta K}
\end{align*}
Por lo que $\pr{r_{12}^2}$ como función de \(T\) será
\begin{equation}
	\pr{r_{12}^2} = \frac{3 k_B}{K} T
\end{equation}
lo cual reproduce el teorema de equipartición.
\qed

\section{Problema 2}
Considere un sistema gaseoso constituido por $N$ moléculas diatómicas no interactuantes, cada una de las cuales posee un momento dipolar eléctrico permanente $\mu$, colocadas en un campo eléctrico externo uniforme de intensidad $E$. La energía de tal molécula está dada por la suma de la energía cinética de traslación, la energía cinética de rotación y la energía potencial de orientación en el campo aplicado:
\begin{equation}
	\epsilon = \frac{p^2}{2 m} + \frac{1}{2 I}\left(p_\theta^2 + \frac{p_\phi^2}{\sin^2\theta}\right) - \mu E \cos\theta
\end{equation}
donde $I$ es el momento de inercia de la molécula.
\begin{enumerate}
	\item[(a)] Estudie la termodinámica de este sistema, incluyendo la polarización eléctrica y la constante dieléctrica. Suponga que el sistema es clásico y que el parámetro $\mu E \ll k_B T$.
\end{enumerate}













\end{document}